\subsection{参数标定与结果对照}
2025 年四大家鱼约恢复 1.8 倍这一事实主要用于约束总量尺度。我们在这一问里有意把标定目标压缩到单一总量事实上。考虑到单一总量事实不足以同时识别多组结构参数，我们固定单区食物网的结构参数，只对增益缩放因子 $gain\_scale$ 做单参数标定，并保持鱼类内部比例不变；$consumer\_scale=0.60$ 仅承担鱼类初值规范化任务。

为了让 2025 年总量约束只作用于禁渔路径，我们让$F_0$ 不进入参数搜索，而在禁渔基线确定后单独构造“不禁渔”情景，同时把湖北段 $\mathrm{CPUE}^*$ 保留为标定后的局地对照。数值求解采用 \texttt{solve\_ivp}，标定结果为 $gain\_scale=1.45$，对应 2025 年四大家鱼总量 1.805 倍、$\mathrm{CPUE}^*$ 1.838 倍，对 1.8 倍与 2.2 倍的相对误差分别为 0.26\% 和 16.47\%。

不禁渔情景由固定反事实参数 $F_0=0.25$ 给出；该值位于预设区间 $[0.10,0.40]$ 中部，只用于比较禁渔与持续捕捞的方向差异。进一步将 $F_0$ 在该区间内等距取 7 个值重复对照，禁渔情景始终优于不禁渔情景（$B_{\text{carp}}$ 差值均为正），且 $F_0\ge 0.15$ 时差值超过 30\%。因此，“禁渔优于不禁渔”的判断不依赖某一个特定的 $F_0$ 取值。

